\section{Introduction}
Personalization has become a key factor in improving user engagement in various digital experiences, including travel and tourism. As travellers seek more relevant and tailored recommendations, AI advancements, particularly large language models (LLM), presents new opportunities to adapt content to align with individual preferences dynamically. One such application is personalized descriptions of Points of Interest (POIs), such as hotels, restaurants, museums, and gardens. Although traditional POI descriptions are often static and generalized, AI-driven personalization can refine these descriptions to match user-specific interests, past experiences, and contextual relevance, enhancing both engagement and comprehension.

Much prior research has focused on POI recommendation and ranking systems, where AI models suggest locations based on user preferences, behaviors, and contextual factors. However, less attention has been given to how the content of each POI itself can be adapted to better engage users. When individuals access POI descriptions, the same information is typically presented to all users, despite variations in their backgrounds, interests, and motivations to visit a place. For example, a museum description could emphasize historical significance for a history enthusiast while focusing on interactive exhibits for a family traveler. Generative AI and LLMs provide a unique opportunity to dynamically tailor content without altering its factual meaning, ensuring that the description remains relevant and appealing to different audiences.

To explore this potential, a prototype system was developed that allows users to enter their preferences and interests and then dynamically presents two versions of each POI description: the original and an LLM-adapted version. Participants were asked to evaluate the descriptions based on alignment with their preferences, perceived appeal, and familiarity with the location. The key objective is to assess whether AI-generated adaptations are acceptable and meaningful, effectively convey the original value of POI, and improve user engagement.

This study contributes to the growing field of AI-powered content personalization in tourism by investigating the role of LLMs in dynamic adaptation of textual content. By analyzing participant feedback, this study aims to provide insight into how AI-generated descriptions can improve engagement while preserving accuracy and authenticity.

% The remainder of this paper is structured as follows. Section 2 details the methodology, including the experimental setup, data collection process, and evaluation metrics. Section 3 presents the results and analysis of user perceptions across significance, trustworthiness, and clarity dimensions. Section 4 discusses the implications of the findings and limitations of the study. Section 5 outlines directions for future research, and Section 6 concludes with key takeaways and practical implications for tourism platforms.
